\subsection{7.18 Complete System Matrix Derivation - Equations (7) and (8)}

This subsection provides the complete step-by-step derivation of the system matrix $[A]$ and state-vector transformation, matching the equations shown in the attached reference material.

\subsubsection{7.18.1 State Vector Definition}

The complete state vector is defined as:

\begin{equation}
\boxed{\begin{bmatrix}
X_{01} \\
Y_{01} \\
X_{21} \\
Y_{21} \\
T_1 \\
T_2
\end{bmatrix} = [A] \begin{bmatrix}
X_{01} \\
Y_{01} \\
X_{21} \\
Y_{21} \\
T_1 \\
T_2
\end{bmatrix}}
\end{equation}

where the notation is:
\begin{itemize}
    \item $X_{01}, Y_{01}$ = End-effector (point M) position in base frame (frame 0)
    \item $X_{21}, Y_{21}$ = End-effector velocity components (frame 2 perspective or Cartesian derivatives)
    \item $T_1, T_2$ = Cable tensions
\end{itemize}

\textbf{Alternative notation clarification:}

The subscript notation can be interpreted as:
\begin{align}
X_{01} &= x_m \text{ (position of point M from origin 0)} \\
Y_{01} &= y_m \text{ (position of point M from origin 0)} \\
X_{21} &= \dot{x}_m \text{ (velocity, time derivative = 1st time derivative)} \\
Y_{21} &= \dot{y}_m \text{ (velocity)} \\
T_1, T_2 &= \text{Cable tensions (state-like variables)}
\end{align}

\subsubsection{7.18.2 Time Derivatives of State Components}

The first and second components represent position, so their time derivatives are velocities:

\begin{equation}
\dot{X}_{01} = X_{21}
\end{equation}

\begin{equation}
\dot{Y}_{01} = Y_{21}
\end{equation}

\textbf{State equations for position (rows 1-2 of [A]):}

\begin{equation}
\dot{X}_{01} = 0 \cdot X_{01} + 0 \cdot Y_{01} + 1 \cdot X_{21} + 0 \cdot Y_{21} + 0 \cdot T_1 + 0 \cdot T_2
\end{equation}

\begin{equation}
\dot{Y}_{01} = 0 \cdot X_{01} + 0 \cdot Y_{01} + 0 \cdot X_{21} + 1 \cdot Y_{21} + 0 \cdot T_1 + 0 \cdot T_2
\end{equation}

\subsubsection{7.18.3 Time Derivatives of Velocity Components}

The velocity components evolve according to the equations of motion from Euler-Lagrange:

\begin{equation}
\mathbf{M}(\mathbf{q}) \ddot{\mathbf{q}} + \mathbf{G}(\mathbf{q}) = \boldsymbol{\tau} + \mathbf{J}_{\rho}^T(\mathbf{q}) \mathbf{T}
\end{equation}

Solving for joint accelerations:

\begin{equation}
\ddot{\mathbf{q}} = \mathbf{M}^{-1} \left[ \mathbf{J}_{\rho}^T(\mathbf{q}) \mathbf{T} - \mathbf{G}(\mathbf{q}) \right]
\end{equation}

\textbf{Transforming to Cartesian acceleration:}

The Cartesian accelerations are related to joint accelerations through:

\begin{equation}
\ddot{\mathbf{r}}_m = \mathbf{J}_q(\mathbf{q}) \ddot{\mathbf{q}} + \dot{\mathbf{J}}_q(\mathbf{q}, \dot{\mathbf{q}}) \dot{\mathbf{q}}
\end{equation}

where $\mathbf{J}_q$ is the joint Jacobian.

Therefore:

\begin{equation}
\begin{bmatrix}
\dot{X}_{21} \\
\dot{Y}_{21}
\end{bmatrix} = \begin{bmatrix}
\ddot{x}_m \\
\ddot{y}_m
\end{bmatrix} = f(X_{01}, Y_{01}, X_{21}, Y_{21}, T_1, T_2)
\end{equation}

\textbf{Detailed expansion:}

The acceleration components depend on:

1. **Configuration effects** (position-dependent): $-\sin q_i$, $\cos q_i$ terms
2. **Centrifugal/Coriolis effects** (velocity-dependent): $(\dot{q}_1)^2$ terms
3. **Gravity terms**: $g \cos(q_i)$ components
4. **Tension coupling**: $\mathbf{J}_{\rho}^T \mathbf{T}$ terms

\subsubsection{7.18.4 Gravity Component Decomposition}

From equation (7) in your attachment, the constant term vector on the RHS shows:

\begin{equation}
\text{Gravity/Inertial Effects} = \begin{bmatrix}
0 \\
m_1 g \\
\frac{1}{2} m_1 g \cos q_1 \\
0 \\
m_2 g \\
\frac{1}{2} m_2 g \cos(q_1 + q_2)
\end{bmatrix}
\end{equation}

These represent:
\begin{itemize}
    \item Row 1: Position equation (no constant effect)
    \item Row 2: Gravity effect on link 1 (vertical component $m_1 g$)
    \item Row 3: Torque component from link 1 mass
    \item Row 4: Velocity equation (no constant effect)
    \item Row 5: Gravity effect on link 2
    \item Row 6: Torque component from link 2 mass
\end{itemize}

\subsubsection{7.18.5 Complete System Matrix [A] - Equation (8)}

The system matrix from equation (8) has the following structure:

\begin{equation}
[A] = \begin{bmatrix}
1 & 0 & 1 & 0 & 0 & 0 \\
0 & 1 & 0 & 1 & 0 & 0 \\
0 & 0 & -1, \sin q_1 & l_1 \cos q_1 & 0 & 0 \\
0 & 0 & -1 & 0 & -\cos \alpha_1 & -\cos \alpha_2 \\
0 & 0 & 0 & -1 & -\sin \alpha_1 & -\sin \alpha_2 \\
0 & 0 & 0 & 0 & l_2 \sin(q_1 + q_2 + \alpha_1) & l_2 \sin(q_1 + q_2 + \alpha_2)
\end{bmatrix}
\end{equation}

where:
\begin{itemize}
    \item $q_1, q_2$ = Joint angles
    \item $\alpha_1, \alpha_2$ = Cable angles at attachment point
    \item $l_1, l_2$ = Link lengths
\end{itemize}

\textbf{Detailed matrix element explanation:}

\textbf{Upper-left $2 \times 2$ block (Position kinematics):}

\begin{equation}
\begin{bmatrix}
\dot{X}_{01} \\
\dot{Y}_{01}
\end{bmatrix} = \begin{bmatrix}
0 & 0 & 1 & 0 & 0 & 0 \\
0 & 0 & 0 & 1 & 0 & 0
\end{bmatrix} \begin{bmatrix}
X_{01} \\
Y_{01} \\
X_{21} \\
Y_{21} \\
T_1 \\
T_2
\end{bmatrix}
\end{equation}

This shows that position derivatives equal the velocity components (identity relationship).

\textbf{Rows 3-4 (Acceleration kinematics):}

\begin{equation}
\begin{bmatrix}
\dot{X}_{21} \\
\dot{Y}_{21}
\end{bmatrix} = \begin{bmatrix}
0 & 0 & -1 \sin q_1 & l_1 \cos q_1 & 0 & 0 \\
0 & 0 & -1 & 0 & -\cos \alpha_1 & -\cos \alpha_2
\end{bmatrix} \begin{bmatrix}
X_{01} \\
Y_{01} \\
X_{21} \\
Y_{21} \\
T_1 \\
T_2
\end{bmatrix}
\end{equation}

Breaking this down:
\begin{itemize}
    \item Column 3 term: $-\sin q_1 \cdot X_{21}$ - represents joint angle dependent velocity coupling
    \item Column 4 term: $l_1 \cos q_1 \cdot Y_{21}$ - link length and angle coupling
    \item Column 5-6 terms: $-\cos \alpha_i \cdot T_i$ - cable tension effects on horizontal acceleration
\end{itemize}

\textbf{Rows 5-6 (Cable tension dynamics):}

\begin{equation}
\begin{bmatrix}
\dot{T}_1 \\
\dot{T}_2
\end{bmatrix} = \begin{bmatrix}
0 & 0 & 0 & -1 & -\sin \alpha_1 & -\sin \alpha_2 \\
0 & 0 & 0 & 0 & l_2 \sin(q_1+q_2+\alpha_1) & l_2 \sin(q_1+q_2+\alpha_2)
\end{bmatrix} \begin{bmatrix}
X_{01} \\
Y_{01} \\
X_{21} \\
Y_{21} \\
T_1 \\
T_2
\end{bmatrix}
\end{equation}

This represents:
\begin{itemize}
    \item Cable tensions evolve based on velocity and configuration
    \item The $\sin \alpha_i$ terms represent vertical components of cable tensions
    \item The $\sin(q_1+q_2+\alpha_i)$ terms represent geometric coupling
\end{itemize}

\subsubsection{7.18.6 Physical Interpretation of Matrix Elements}

\textbf{Element $A_{ij}$ meaning:}

Each element $A_{ij}$ represents the influence of state variable $j$ on the time derivative of state variable $i$.

\begin{equation}
\dot{x}_i = \sum_{j=1}^{6} A_{ij} x_j
\end{equation}

\textbf{Key elements and their meanings:}

\begin{enumerate}
    \item $A_{1,3} = 1$: Position rate = velocity (kinematic constraint)
    
    \item $A_{2,4} = 1$: Vertical position rate = vertical velocity
    
    \item $A_{3,3} = -\sin q_1$: Velocity coupling through joint angle
    
    \item $A_{3,4} = l_1 \cos q_1$: Link length effect on acceleration
    
    \item $A_{3,5} = 0$, $A_{3,6} = 0$: No direct tension effect on horizontal acceleration (in linear approximation)
    
    \item $A_{4,5} = -\cos \alpha_1$: Cable 1 tension affects vertical acceleration
    
    \item $A_{4,6} = -\cos \alpha_2$: Cable 2 tension affects vertical acceleration
    
    \item $A_{5,4} = -1$: Vertical velocity affects tension rate
    
    \item $A_{5,5} = -\sin \alpha_1$: Cable angle effects
    
    \item $A_{6,6} = l_2 \sin(q_1+q_2+\alpha_2)$: Geometric coupling with link 2
\end{enumerate}

\subsubsection{7.18.7 Deriving Matrix Elements from First Principles}

\textbf{Forward Kinematics (Rows 1-2):}

From:
\begin{equation}
X_{01} = l_1 \cos q_1 + l_2 \cos(q_1+q_2)
\end{equation}

\begin{equation}
Y_{01} = l_1 \sin q_1 + l_2 \sin(q_1+q_2)
\end{equation}

Taking time derivatives:

\begin{equation}
\dot{X}_{01} = -l_1 \sin q_1 \cdot \dot{q}_1 - l_2 \sin(q_1+q_2) \cdot (\dot{q}_1 + \dot{q}_2)
\end{equation}

In terms of state variables where $X_{21} \approx \dot{X}_{01}$ (for infinitesimal dt):

\begin{equation}
\dot{X}_{01} = X_{21}
\end{equation}

\textbf{This gives:} $A_{1,3} = 1$, all other $A_{1,j} = 0$.

\textbf{Joint Acceleration (Rows 3-4):}

From the Euler-Lagrange dynamics:

\begin{equation}
\ddot{X}_{21} = \mathbf{J}_{q1,1} \ddot{q}_1 + \mathbf{J}_{q1,2} \ddot{q}_2 + \dot{\mathbf{J}}_{q1} \mathbf{v}
\end{equation}

Inverting the mass matrix and substituting:

\begin{equation}
\ddot{X}_{21} = \sum_j M^{-1}_{11} [J_{\rho,j} T_j - G_j] + \text{velocity-dependent terms}
\end{equation}

The coefficients involving $\sin q_1$, $\cos q_1$, $\alpha_i$ come from:
\begin{itemize}
    \item Jacobian elements (kinematic structure)
    \item Inverse mass matrix (inertial effects)
    \item Cable geometry terms
\end{itemize}

\textbf{Cable Tension Dynamics (Rows 5-6):}

For cable tensions with actuator dynamics:

\begin{equation}
\dot{T}_i = -k_i T_i + k_i \tau_{c,i} + \text{constraint forces}
\end{equation}

The constraint forces depend on:
\begin{itemize}
    \item Cable length rate: $\dot{\rho}_i = \mathbf{J}_{\rho,i} \cdot \mathbf{v}$
    \item Cable tension equilibrium: $\sum T_i \cdot \text{(unit vectors)} = \mathbf{F}_{ext}$
\end{itemize}

This gives the coupling terms involving $\sin \alpha_i$ and geometric factors.

\subsubsection{7.18.8 Special Cases and Simplifications}

\textbf{Case 1: Small angle approximation} ($\sin q_i \approx q_i$, $\cos q_i \approx 1$)

\begin{equation}
[A]_{\text{linear}} \approx \begin{bmatrix}
0 & 0 & 1 & 0 & 0 & 0 \\
0 & 0 & 0 & 1 & 0 & 0 \\
0 & 0 & -q_1 & l_1 & 0 & 0 \\
0 & 0 & -1 & 0 & -\cos \alpha_1 & -\cos \alpha_2 \\
0 & 0 & 0 & -1 & -\sin \alpha_1 & -\sin \alpha_2 \\
0 & 0 & 0 & 0 & l_2 \sin(\alpha_1) & l_2 \sin(\alpha_2)
\end{bmatrix}
\end{equation}

\textbf{Case 2: Horizontal cables} ($\alpha_1 = \alpha_2 = 0$)

\begin{equation}
[A]_{\text{horiz}} = \begin{bmatrix}
0 & 0 & 1 & 0 & 0 & 0 \\
0 & 0 & 0 & 1 & 0 & 0 \\
0 & 0 & -\sin q_1 & l_1 \cos q_1 & 0 & 0 \\
0 & 0 & -1 & 0 & -1 & -1 \\
0 & 0 & 0 & -1 & 0 & 0 \\
0 & 0 & 0 & 0 & l_2 \sin(q_1+q_2) & l_2 \sin(q_1+q_2)
\end{bmatrix}
\end{equation}

\textbf{Case 3: Vertical cables} ($\alpha_1 = \alpha_2 = \pi/2$)

\begin{equation}
[A]_{\text{vert}} = \begin{bmatrix}
0 & 0 & 1 & 0 & 0 & 0 \\
0 & 0 & 0 & 1 & 0 & 0 \\
0 & 0 & -\sin q_1 & l_1 \cos q_1 & 0 & 0 \\
0 & 0 & -1 & 0 & 0 & 0 \\
0 & 0 & 0 & -1 & -1 & -1 \\
0 & 0 & 0 & 0 & l_2 \sin(q_1+q_2+\pi/2) & l_2 \sin(q_1+q_2+\pi/2)
\end{bmatrix}
\end{equation}

\subsubsection{7.18.9 Numerical Computation of [A]}

\textbf{Algorithm:}

\begin{algorithm}[H]
\caption{Compute System Matrix [A] at Configuration $\mathbf{q}$}
\begin{algorithmic}[1]
\REQUIRE Joint angles $q_1, q_2$, Link parameters $l_1, l_2$, Cable angles $\alpha_1, \alpha_2$
\STATE Compute joint Jacobian: $\mathbf{J}_q(\mathbf{q})$
\STATE Compute inverse mass matrix: $\mathbf{M}^{-1}(\mathbf{q})$
\STATE Compute gravity gradient: $\nabla \mathbf{G}(\mathbf{q})$
\STATE Compute cable geometry Jacobian: $\mathbf{J}_{\rho}(\mathbf{q})$
\STATE Assemble upper blocks (kinematics): rows 1-2, rows 3-4
\STATE Assemble tension coupling: $-\mathbf{M}^{-1} \mathbf{J}_{\rho}^T$
\STATE Assemble gravity effects: $-\mathbf{M}^{-1} \nabla \mathbf{G}$
\STATE Combine into $6 \times 6$ matrix $[A]$
\RETURN $[A](\mathbf{q})$
\end{algorithmic}
\end{algorithm}

\subsubsection{7.18.10 Connection to Singularity Analysis}

The system matrix $[A]$ reveals critical information about singularities:

\begin{enumerate}
    \item \textbf{Rank deficiency}: At singular configurations, rows become linearly dependent
    
    \item \textbf{Large eigenvalues}: Near singularity, $\lambda_{\max}([A]) \to \infty$
    
    \item \textbf{Instability indicator}: Eigenvalues with large positive real parts indicate unstable modes
    
    \item \textbf{Controllability loss}: At singularity, some state variables become uncontrollable
    
    \begin{equation}
    \text{rank}[B, AB, A^2B, \ldots] < n \text{ at singularity}
    \end{equation}
\end{enumerate}

\textbf{Monitoring singularity proximity:}

\begin{equation}
\text{Condition number: } \kappa([A]) = \frac{\lambda_{\max}([A])}{\lambda_{\min}([A])} \gg 1 \text{ near singularity}
\end{equation}

\subsubsection{7.18.11 Linearized vs. Nonlinear System}

\textbf{Linearized around operating point:}

\begin{equation}
\dot{\mathbf{x}} = [A](\mathbf{q}_0) \mathbf{x} + \mathbf{B}(\mathbf{q}_0) \mathbf{u}
\end{equation}

Valid for small deviations: $\|\mathbf{x} - \mathbf{x}_0\| < \epsilon$

\textbf{Nonlinear system:}

\begin{equation}
\dot{\mathbf{x}} = \mathbf{f}(\mathbf{x}, \mathbf{u})
\end{equation}

where $[A] = \frac{\partial \mathbf{f}}{\partial \mathbf{x}}\big|_{\mathbf{x}_0}$

\textbf{Linearization error depends on:}
\begin{itemize}
    \item Configuration ($q_1$, $q_2$ terms are nonlinear)
    \item Velocity (centrifugal terms)
    \item Tension magnitude (coupling strength)
\end{itemize}

For larger deviations, the full nonlinear model must be used for accurate simulation.

\subsubsection{7.18.12 Summary: System Matrix Derivation}

The complete system matrix $[A]$ is derived from:

\begin{enumerate}
    \item \textbf{Kinematics}: Position-velocity relationship (rows 1-2)
    
    \item \textbf{Dynamics}: Joint accelerations from Euler-Lagrange (rows 3-4)
    
    \item \textbf{Actuation}: Cable tension rate equations (rows 5-6)
    
    \item \textbf{Constraints}: Cable length and force balance relationships
    
    \item \textbf{Configuration}: Joint angles $q_1, q_2$ appear explicitly
    
    \item \textbf{Geometry}: Link lengths $l_1, l_2$ and cable angles $\alpha_1, \alpha_2$
\end{enumerate}

The resulting $6 \times 6$ matrix captures:
- ✅ Full nonlinear dynamics in linearized form
- ✅ Coupling between position, velocity, and tension
- ✅ Configuration-dependent effects
- ✅ Singularity sensitivity through matrix conditioning
- ✅ Foundation for control design and stability analysis
